\section*{Sets and Indices}

\begin{itemize}
    \item $T=\{0,1,\dots,5\}$: set of 4-hour time windows, in the order given by \texttt{table\_1.csv}.
    \item $j \in T$: index of a coverage window.
    \item $i \in T$: index of the first window of an 8-hour tour.
    \item For each $i \in T$, let the candidate tour cover windows $i$ and $(i+1)\bmod |T|$.
    \item $A \subseteq T$: set of admissible tour starts,
    \[
        A=\{\, i \in T : \texttt{peak\_indicator}_i \neq \texttt{peak\_indicator}_{(i+1)\bmod |T|}\,\},
    \]
    i.e., only adjacent peak/off-peak pairs are allowed.
\end{itemize}

\section*{Parameters}

\begin{itemize}
    \item $d_j$ (\texttt{min\_waiters[j]}): minimum required number of waiters in window $j$.
    \item $\pi_j$ (\texttt{peak\_indicator[j]}): peak indicator for window $j$.
    \item $q_j$ (\texttt{quality\_weight[j]}): service quality weight for staffing in window $j$.
    \item $c$ (\texttt{waiter\_cost\_per\_shift}): cost of one 8-hour waiter tour.
    \item $B$ (\texttt{waiter\_budget}): total staffing budget.
    \item $a_{ij}$ (\texttt{coverage[(i,0,j)]}): coverage indicator, defined for admissible $i \in A$ by
    \[
        a_{ij}=
        \begin{cases}
            1, & \text{if } j=i \text{ or } j=(i+1)\bmod |T|,\\
            0, & \text{otherwise}.
        \end{cases}
    \]
    Since the code creates at most one pattern per start $i$, the pattern index is fixed to $0$ whenever it exists.
\end{itemize}

\section*{Decision Variables}

\begin{itemize}
    \item $y_i$ (code: \texttt{y[(i,0)]}) for each $i \in A$: integer number of waiters assigned to the admissible 8-hour tour starting in window $i$,
    \[
        y_i \in \mathbb{Z}_{\ge 0}.
    \]
    \item $w_j$ (code: \texttt{waiters\_in\_period[j]}$): implied total number of waiters covering window $j$,
    \[
        w_j=\sum_{i \in A} a_{ij} y_i \qquad \forall j \in T.
    \]
    This is a linear expression in the implementation, not a separate optimization variable.
\end{itemize}

\section*{Objective}

Maximize the total service quality score:
\[
\max \sum_{j \in T} q_j\, w_j
=
\max \sum_{j \in T} q_j \left(\sum_{i \in A} a_{ij} y_i\right).
\]

\section*{Constraints}

\begin{align*}
w_j &= \sum_{i \in A} a_{ij} y_i && \forall j \in T,\\
w_j &\ge d_j && \forall j \in T,\\
c \sum_{i \in A} y_i &\le B,\\
y_i &\in \mathbb{Z}_{\ge 0} && \forall i \in A.
\end{align*}

Equivalently, eliminating the implied expressions $w_j$:
\begin{align*}
\sum_{i \in A} a_{ij} y_i &\ge d_j && \forall j \in T,\\
c \sum_{i \in A} y_i &\le B,\\
y_i &\in \mathbb{Z}_{\ge 0} && \forall i \in A.
\end{align*}

\section*{Notes on Implementation}

\begin{itemize}
    \item The implemented model is a deterministic single-stage integer linear program solved with \texttt{GUROBI\_CMD}. Although the business brief refers to multiple cases, the current code does not introduce scenario indices, nonanticipativity constraints, or scenario-dependent recourse variables.
    \item A feasible 8-hour tour is represented as exactly two adjacent 4-hour windows. The code permits such a tour only when one window is peak and the other is off-peak, i.e., when \texttt{peak\_indicator[i] != peak\_indicator[(i+1)\%6]}.
    \item The implementation creates at most one pattern for each start window $i$; hence the only pattern identifier used in code is \texttt{p = 0}.
    \item The objective maximizes \texttt{service\_quality\_score} as the weighted sum of staffed waiters by window. It does \emph{not} minimize the number of waiters.
    \item Budget is enforced on the number of assigned 8-hour tours via \texttt{waiter\_cost\_per\_shift}. Because all tour costs are identical, this limits the total number of assigned shifts.
    \item If some window had no admissible covering tours, then in the code its staffing expression would be the empty sum (equal to $0$), and the corresponding minimum-coverage constraint would make the model infeasible. In this instance, admissible tours are generated from the data-driven peak/off-peak pattern.
\end{itemize}
